% ============================================================
%  §6  Dynamics
% ============================================================

\section{Dynamics}
\label{sec:dynamics}

\subsection{Sedimentation}

Sedimentation is the process by which repeated token--field interaction
deepens potential wells and stabilizes conceptual attractors.
It proceeds in three stages:

\textbf{Stage 1: Initial perturbation.}
A novel input $\Delta$ generates a $\sigma_{\mathrm{drive}}$ spike,
locally increasing $I$ and creating a transient potential gradient.
The information flow $\JI$ responds by redirecting toward the
perturbed region.

\textbf{Stage 2: Reinforcement.}
If the perturbation is repeated --- through subsequent exposure to
similar input or through internal recursive activation --- the local
gradient in $\Phi$ is reinforced.
Token--field coupling \eqref{eq:coupling} ensures that each reuse
deepens the associated well.

\textbf{Stage 3: Stabilization.}
As the well deepens, the traversal cost of the corresponding conceptual
region decreases.
New trajectories are drawn toward the attractor, further reinforcing it.
The concept becomes self-sustaining.

\subsection{Forgetting}

Forgetting corresponds to the relaxation of $\Phi$ toward a flatter
configuration when reuse frequency drops:
\begin{equation}
  \sigma_{\mathrm{dissip}} = \Gamma C.
\end{equation}
This is not erasure of stored symbols but a smooth redistribution of
stability across the field.
Forgotten concepts do not disappear; their wells shallow, raising the
traversal cost until they are effectively inaccessible.
Reactivation (``remembering'') corresponds to re-deepening a shallow
well through renewed token--field interaction.

\subsection{Conceptual Change}

Conceptual change --- the shift in meaning, scope, or application of
an established concept --- corresponds to a topological change in
$\Phistar$: the merger of two wells, the splitting of one well into
two, or the displacement of a well from one region of the landscape
to another.

These changes require sustained perturbation sufficient to overcome
the depth of existing attractors.
The effective barrier term $\sigma_{\mathrm{barrier}}$ plays a
protective role here: as $C\to\Cmax$ in a deeply entrenched concept,
the projected resistance grows, making radical conceptual change
increasingly costly.
This should be understood as an effective description of overclosure
within the conceptual observation window, not as an externally imposed
wall.
This accounts for the resistance of deeply held beliefs and intuitions
to revision.

\subsection{Insight as Rapid Topology Change}

The phenomenology of insight --- the sudden sense that a problem has
been resolved or that disparate ideas have been connected --- corresponds
to a rapid change in the local topology of $\Phi$.
In field-theoretic terms, this is a bifurcation event: two previously
separated wells merge, or a barrier that had prevented flow between
two regions suddenly collapses.

The subjective sense of suddenness reflects the threshold character
of topological transitions: small perturbations accumulate below the
threshold of perceptual change, then produce a discrete topological
event when the threshold is crossed.
